\documentclass[../main.tex]{subfiles}
\begin{document}
%==========================================TIÊU ĐỀ TẬP========================================
\begin{table}[h]
    \begin{adjustwidth}{-1cm}{}
        \begin{tabular}{>{\centering\arraybackslash}p{6cm}|>{\raggedright\arraybackslash}p{12.5cm}}
            \multirow{5}{6cm}
            {\\[-40pt]
                \flaregothic\fontsize{20pt}{20pt}\selectfont \centering
                \color{eptype}HISTOIRE\color{black}\\
                \freshbot\fontsize{72pt}{72pt}\selectfont
                \color{epnum}20\\[6pt]
                \cafeteria\fontsize{20pt}{20pt}\selectfont\color{epnum}(D-20)\color{black}
                \\[18pt]
                \color{black}
            }
             &
            {\fotskip\fontsize{18pt}{18pt}\selectfont \ruby{史}{し}\ruby{上}{じょう}\ruby{最}{さい}\ruby{大}{だい}の\ruby{口}{こう}\ruby{論}{ろん}}
            \\[6pt]
             & {\cafeteria\fontsize{18pt}{18pt}\selectfont The Great Verbal Showdown} \\[4pt]
             & {\vnmsans\fontsize{18pt}{18pt}\selectfont Trận đấu khẩu lịch sử} \\[8pt]
             & {\caviar{\fontsize{14pt}{14pt}\selectfont Nhân vật: \emp{suzuno2} \emp{kaigou2} \emp{ryuuhou2} \emp{kuon2} \emp{game} \emp{kson} \emp{delu} \emp{naomi2}}} \\[8pt]
        \end{tabular}
    \end{adjustwidth}
\end{table}
%=========================================TRUYỆN CHÍNH========================================
\color{black}

\txtbx[2]{
    {\color{vol} \uline{Tóm tắt:}} Tối thứ Năm, Ryuuhou quên rửa bát lần thứ ba trong tuần. Lời nhắc bình thường của Kaigou nhanh chóng biến thành cuộc tranh luận về việc ai làm nhiều việc nhà hơn. Kson và Delutaya theo dõi như đang xem chương trình trực tiếp. Kuon về nhà, nghe hết hai phía rồi buộc cả hai cùng dọn bếp và viết lại cách bàn giao việc. Trong lúc rửa bát, một câu đùa ngớ ngẩn của Ryuuhou khiến Kaigou bật cười, khép lại trận cãi vã.
}

\pov{Hikawa Suzuno}

\lang{Etsunan}

Ngày hôm đó, chị Kaigou và Ryuuhou suýt nữa thì biến cả tầng năm thành một phiên tòa xử án, một chiến trường kịch tính, rồi cuối cùng lại buộc phải dồn hết sức lực cho một buổi tổng vệ sinh không ai dự đoán trước.

Tôi biết rõ câu chuyện ấy bắt nguồn từ đâu.

Tất cả chỉ bắt đầu từ một chiếc bát.

Nói chính xác hơn, nó khởi nguồn từ chiếc bát trắng bạc còn nằm quên trong bồn rửa ngay sau khi bữa tối vừa kết thúc.

Nhưng nếu chỉ có vậy, mọi sự việc chắc chắn đã dừng lại sau một câu nhắc nhở nhẹ nhàng, chẳng có gì to tát xảy ra.

Vấn đề nằm ở chỗ, chiếc bát đó xuất hiện vào tối thứ Năm,tức là lần thứ ba trong cùng một tuần mà Ryuuhou quên không rửa bát của mình, hai lần trước đó đã bị bỏ qua không ai nhắc nhở.

Vấn đề thứ hai lại càng lớn hơn: chị tôi đã ghi lại cả hai lần quên đó cẩn thận vào một cuốn sổ nhỏ.

Còn Ryuuhou, em ấy cũng chuẩn bị sẵn một danh sách khác, dài không kém, để đối phó.

Tôi chỉ thực sự hiểu ra sự nghiêm trọng của vấn đề sau khi cuộc tranh luận nảy lửa chính thức bùng nổ, không thể dập tắt được.

Lúc ấy, tôi còn nghĩ mình chỉ cần nhắc hai người bình tĩnh rồi tiếp tục dọn dẹp như mọi tối. Tôi không ngờ chiếc bát nhỏ kia lại chạm đúng vào những điều cả hai đã giữ trong lòng.

\il

Tối hôm đó, tôi đang cầm miếng vải lau, chăm chỉ lau sạch mâm cơm, nơi mọi người thường tập trung sau bữa ăn.

Naomi ngồi trên chiếc ghế bành đệm mềm góc phòng, đôi tay nhỏ bé cẩn thận xếp từng chiếc khăn ăn thành những chồng nhỏ gọn gàng, rất đều đặn.

Anh Kuon vẫn chưa trở về từ công ty, văn phòng của anh ấy có việc nên phải ở lại muộn hơn thường ngày.

Chiếc đồng hồ thông minh của chị tôi mà Game trú ngụ được đặt ở giữa bàn, màn hình màu đen tắt lịm, nhưng tôi biết rõ ông ấy vẫn đang lắng nghe mọi chuyện diễn ra trong căn nhà, chẳng có gì qua được mắt ông ấy.

Hơn nửa tiếng đồng hồ đã trôi qua kể từ khi bữa tối chính thức kết thúc, những đĩa lớn đựng thức ăn đã được dọn bớt, mùi thơm của các món ăn vừa rồi vẫn còn vương lại nhẹ nhàng trong không khí phòng ăn. Mọi người đã vừa xong bữa, nhưng trong bếp vẫn còn ngổn ngang vài công việc nhỏ nhặt chưa ai hoàn thành: mấy chiếc nồi lớn còn bám lại bọt dầu sau khi nấu, vài chiếc bát đĩa dơ nằm rải rác trên quầy bếp, chưa ai gom lại bồn rửa để dọn dẹp hết.

Chị tôi đang đứng ở bồn rửa, cẩn thận rửa những chiếc nồi lớn vừa dùng để nấu ăn, nước chảy róc rách từ vòi.

Còn Ryuuhou thì ngồi dựa lưng vào ghế gỗ, mắt dán vào tờ giấy ghi chú, đang xem lại lịch trình cho buổi tập của ban nhạc mà em ấy tham gia vào ngày mai, tập trung hết sức.

Trên bàn, một chiếc bát, một đôi đũa gỗ và một chiếc đĩa nhỏ còn nằm nguyên vị, chưa ai động vào. Đó là phần đĩa của Ryuuhou mà em ấy để lại sau khi ăn xong.

Hai lần liền, chị tôi liếc mắt qua chồng bát đĩa còn bày nguyên trên bàn, mỗi lần lại lắc đầu nhẹ nhàng rồi phất ra một tiếng thở dài đầy bực bội. 

Khi sự phiền muộn dần tích tụ, đến lần thứ ba chị không còn chịu được nữa, bàn tay vẫn đang cọ rửa chiếc nồi lớn bỗng dừng lại, chị nhanh chóng vắt nhanh giọt nước còn đọng trên ngón tay vào chiếc khăn vải treo cạnh bồn rửa, tay áo ướt một mảng cũng không còn bận tâm.

``Ryuu, em ăn xong lâu rồi đúng không?'' chị cố giữ giọng bình tĩnh, gọi với ra từ bếp về phía góc bàn nơi Ryuuhou đang ngồi.

Nhưng em gái vẫn chăm chú dán mắt vào tờ giấy lịch tập ban nhạc trải trước mặt, ngón tay còn chỉ vào một dòng chữ nhỏ, thậm chí không ngước lên lấy một cái: ``Rồi ạ, em ăn xong từ hồi nào rồi.''

``Vậy sao bát đĩa của em vẫn cứ nằm nguyên đó, không chịu đem đi rửa cho gọn?'' giọng chị Kaigou bắt đầu có chút không vui, sự phiền muộn đã len lỏi vào câu chữ.

``Em định lát nữa rồi rửa mà,'' Ryuuhou đáp gọn lỏn, vẫn không ngẩng đầu, dường như lịch tập quan trọng hơn mọi việc xung quanh.

Lúc này Kaigou không còn nhịn được nữa, chị vặn chặt vòi nước để dòng nước róc rách lập tức dừng lại, lau khô hẳn hai bàn tay ướt sũng vào khăn, giọng đầy chất vấn: ``Lát nữa là bao giờ? Lúc nào em mới gọi là `lát nữa'?''

``Thì lúc em xem xong hết đống lịch tập này, sắp xếp xong thời gian cho tuần tới thôi.'' Ryuuhou chỉ vào tờ giấy, có vẻ vô cùng bận rộn.

``Câu đó em đã nói y hệt vào hôm thứ Hai đầu tuần rồi,'' chị nhắc lại, giọng đã có chút gắt gỏng, sự bực bội không còn che giấu được. ``Lúc đó em cũng nói `lát nữa rửa', cuối cùng thì bát đĩa đó để nguyên trong bồn cho đến sáng hôm sau.''

Lần này Ryuuhou mới ngẩng đầu lên, mặt có chút ngạc nhiên như không nhớ lại chuyện đó, phản ứng ngay: ``Hôm thứ Hai em đâu có quên? Em đã rửa hết cả nồi lớn lẫn cái chảo chiên đó mà, bếp lúc ấy sạch bóng nhẵn rồi!''

``Được, vậy còn hôm thứ Tư thì sao?'' chị tiếp tục chất vấn, không chịu thua. ``Hôm đó em lại quên rửa bát của mình, đúng không? Lần nào cũng có lý do.''

``Hôm thứ Tư em đã dọn sạch bàn ăn rồi mà,'' Ryuuhou ngay lập tức phản bác, có vẻ oan ức. ``Em đã xếp gọn mọi thứ, lau sạch mặt bàn, đâu có bỏ bừa bộn đâu chị?''

``Em đẩy toàn bộ đống bát đĩa dơ vào bồn rửa rồi chuồn thôi, đâu có tự tay rửa chúng đi,'' Kaigou nói thẳng, không muốn quanh co gì nữa. ``Đấy mà gọi là dọn dẹp à?''

``Nhưng ít nhất em cũng đã dọn bàn ăn, mặt bàn sạch sẽ mà!'' Ryuuhou vẫn cố gắng cãi lại, không chịu nhận lỗi.

Kaigou nhìn thẳng vào em gái mình một lúc lâu, lặng thinh không nói gì, cố gắng hít sâu một hơi để giữ bình tĩnh, không muốn chuyện này trở thành cuộc cãi vã to tiếng giữa hai người. Giọng chị vẫn cố gắng đều đặn, không muốn mọi chuyện đi quá xa: ``Ryuu, chị thật sự không muốn cãi nhau với em. Chỉ là em làm ơn mang bát đĩa của mình đi rửa trước khi lên phòng ngủ được không? Chị không thể chịu được cảnh đống bát đĩa bẩn đó để qua đêm trong bếp, bồn rửa đã đầy rồi.''

Nghe vậy, Ryuuhou thở một hơi dài như thể chịu rất nhiều thiệt thòi, gấp vội tờ lịch lại trên bàn, giọng có chút bực bội: ``Được rồi được rồi, em đi rửa. Chị cứ nhắc đi nhắc lại mãi câu đó, nghe như thể em chưa bao giờ làm bất cứ việc gì trong nhà cả, toàn ăn bám chị thôi.''

Bàn tay Kaigou vốn đang định vươn ra lấy chiếc bát trên bàn, bỗng dưng khựng lại giữa không trung, bất ngờ trước phản ứng của em gái. Cả không khí trong phòng bỗng thay đổi, căng thẳng dần hiện rõ.

Tôi cũng lập tức dừng việc lau bàn, tay cầm miếng vải vẫn đang chạm vào mặt bàn, lặng lẽ cảm nhận sự ngột ngạt đang lan tỏa khắp căn phòng. Mọi người đều im lặng, không ai dám nói gì. Chiếc bát trắng mà Ryuuhou vừa định đứng lên cầm lấy vẫn nằm yên nguyên vị trên bàn, chưa có một ai nhấc nó đi.

Tôi muốn lên tiếng bảo hai chị em đừng cãi nhau vì một chiếc bát. Nhưng nhìn vẻ mặt của họ, tôi hiểu rằng họ không còn tranh luận về chiếc bát nữa. Nếu chỉ khuyên bỏ qua, những điều chưa được nói ra có lẽ vẫn sẽ nằm lại đó.

``Chị đâu có bao giờ nói em chưa từng đóng góp gì cho việc nhà cả,'' Kaigou nhấn mạnh, giọng vẫn còn giữ được bình tĩnh nhưng đã bắt đầu có chút bực bội dâng lên.

``Nhưng mà trong mắt chị, có phải lúc nào cũng chỉ có mỗi lần em quên làm việc thôi sao? Không bao giờ chị nhớ đến những lúc em tự giác làm thêm thứ gì đó sao?'' Ryuuhou đáp lời, hai má em hơi ửng đỏ vì cảm thấy bị đối xử bất công.

``Đó là vì chị đang nói về đúng việc em vừa quên hôm nay chứ có phải lôi hết tất cả mọi chuyện ra đâu.''

``Em biết mà. Nhưng cái cách chị nhắc nhở cứ như thể em là kẻ vô tích sự, chẳng bao giờ nhấc ngón tay giúp đỡ gì trong căn nhà này cả!''

Kaigou khoanh đôi tay trước ngực, thở một hơi dài để cố gắng giữ bình tĩnh trước những lời phàn nàn của em gái: ``Chị có bao giờ nói ra những lời đó đâu mà em lại nghĩ vậy?''

``Chị không cần phải nói ra lời đâu. Chỉ cần nghe cách chị nhắc nhở em là em đã cảm nhận được rồi,'' Ryuuhou nhấn mạnh, em không chỉ dựa vào lời nói mà còn dựa vào cả thái độ của chị khi nói chuyện.

``Vậy thì em muốn chị phải nói với em thế nào mới đúng ý? Cái cách nào mới có thể khiến em cảm thấy thoải mái khi được nhắc nhở?'' Kaigou hỏi lại, thực sự bối rối không biết phải làm sao để giải quyết mâu thuẫn này.

``Em cũng không biết nữa\dots\ Có lẽ chị đừng ghi nhớ mọi lỗi nhỏ em mắc phải như thể đang lập một hồ sơ tội phạm cho em vậy,'' Ryuuhou đưa ra ý kiến, em cảm thấy bị theo dõi và ghi nhận mọi lỗi lầm của mình một cách quá mức.

Nghe những lời đó, chị tôi lập tức quay sang ngăn kéo nhỏ nằm bên cạnh tủ bếp, nơi chị thường để những đồ dùng cá nhân. Chị kéo nhẹ ngăn kéo ra, ròi cầm lên một cuốn sổ nhỏ có bìa màu xanh đậm trông rất gọn gàng.

Tôi đứng ở đó, có thể nhìn rõ nhãn dán được dán ngay trên bìa cuốn sổ, chữ viết tay cẩn thận ghi rõ: ``Lịch việc nhà - tháng Năm''.

Ryuuhou chớp mắt vài lần, có vẻ bất ngờ không ngờ chị lại thực sự có một cuốn sổ để ghi chép những việc này. Chị tôi sau đó lật cuốn sổ đến trang đã được đánh dấu sẵn bằng một mẩu giấy nhỏ, mở ra cho em gái có thể nhìn thấy rõ những dòng chữ bên trong.

``Chị không hề lập hồ sơ phạm tội gì cả. Đây chỉ là lịch trực nhật việc nhà thôi,'' Kaigou giải thích, muốn cho em thấy rằng mình không có ý định xấu gì.

``Thế thì tại sao chị lại phải ghi cả ngày và giờ cụ thể cho mỗi việc? Chẳng lẽ không thể nhớ được sao?'' Ryuuhou hỏi tiếp, vẫn chưa thể hiểu được lý do chị phải ghi chép chi tiết đến vậy.

``Để chị có thể nhớ rõ ai đã nhận phần việc nào, tránh tình trạng sót việc hay hiểu lầm về trách nhiệm của từng người.'' 

``Ở đây còn có cả dấu sao đỏ nữa, chị ghi dấu này để làm gì?'' Ryuuhou chỉ vào một chi tiết nhỏ trên trang giấy, thắc mắc về ý nghĩa của những dấu hiệu đặc biệt đó.

``Đó là những ngày chị đã phải nhắc em lần thứ hai để em hoàn thành việc nhà của mình.'' 

Em gái chống cả hai tay lên mặt bàn, vẻ mặt khó tin trước những gì chị vừa nói. ``Nghe còn tệ hơn cả hồ sơ phạm tội nữa đó chị ơi.''

Ở góc phòng, Naomi đang xếp khăn ăn cũng phải dừng lại động tác của mình. Em nhìn qua Ryuuhou, rồi lại nhìn sang Kaigou, đôi mắt to tròn đầy sự bối rối.

Tôi biết trong lòng mình rằng tôi nên nói gì đó để can thiệp, xoa dịu bầu không khí đang ngày càng căng thẳng này. Nhưng đồng thời tôi cũng hiểu rõ rằng nếu mình chen ngang vào không đúng lúc, hai chị em chắc chắn sẽ cùng quay sang hỏi ý kiến của tôi, muốn tôi đứng ra làm trọng tài.

Và lúc này đây, tôi vẫn chưa tìm ra được câu trả lời công bằng nào cho câu chuyện phức tạp giữa hai người.

Chị Kaigou mở rộng cuốn sổ ra hơn nữa, xoay nó hoàn toàn về phía Ryuuhou để em có thể đọc rõ từng dòng chữ chị đã ghi. ``Thứ Hai tuần này: em nhận trách nhiệm rửa bát sau bữa tối. Nhưng sáng hôm sau chị vẫn thấy chiếc bát đó còn nằm nguyên trong bồn rửa, không ai động đến.'' 

Ryuuhou nghiêng đầu, chăm chú nhìn vào trang giấy trước mặt như đang cố gắng nhớ lại chi tiết sự việc hôm đó: ``Nhưng hôm đó em đã rửa cái nồi lớn mà chị dùng để nấu ăn mà!''

``Cái nồi đó lẽ ra là phần việc của chị hôm đó, không phải của em.''

``Nhưng hôm đó chị bảo em rửa giúp chị vì chị đang bận gọi điện thoại nói chuyện với mẹ mà,'' Ryuuhou bảo vệ quan điểm của mình, cho rằng mình đã làm đúng những gì chị yêu cầu.

``Chị chỉ nhờ em tráng qua cái nồi thôi, chứ không phải em thay thế hoàn toàn toàn bộ phần việc của chị bằng phần việc của em!'' chị Kaigou cau mày, nhấn mạnh rằng sự hiểu lầm đã bắt đầu từ đó.

``Chị đâu có nói rõ ràng như thế lúc đó đâu. Nếu chị nói rõ ra thì em đã hiểu ngay rồi!'' cô en gái phản bác lại, cho rằng lỗi nằm ở việc chị không diễn đạt rõ ràng yêu cầu của mình.

``Sao em không hỏi lại chị một câu để chắc chắn về yêu cầu đó? Tại sao phải tự mình hiểu lầm như vậy?'' Kaigou hỏi, cảm thấy bức xúc vì sự việc đơn giản lại trở nên phức tạp chỉ vì thiếu sự trao đổi.

``Còn chị thì sao? Sao chị không nói rõ ràng hơn từ đầu để em không thể hiểu lầm được?'' Ryuuhou đáp lại câu hỏi của chị bằng chính câu hỏi đó, đổ lại trách nhiệm cho Kaigou.

Hai người nhìn chằm chằm vào mắt nhau, không ai chịu nhường bước một chút nào. Cả hai đều tin rằng mình đúng, và không ai muốn là người đầu tiên thừa nhận sai lầm.

Đúng lúc đó, chiếc đồng hồ thông minh trên bàn có Game đang trú ngụ bật sáng màn hình lên, phá vỡ sự căng thẳng giữa hai người.

``Tôi có thể truy xuất lại toàn bộ lịch sử phân công việc nhà trong nhóm tin nhắn của gia đình để làm rõ sự việc. Mọi dữ liệu đều được lưu trữ đầy đủ,'' ông ấy đưa ra đề xuất, với khả năng truy xuất dữ liệu của mình, ông có thể giúp làm rõ mọi hiểu lầm.

Kaigou lập tức chỉ vào chiếc đồng hồ, không chần chừ gì nữa. ``Được thôi. Ông mở nó ra đi, để mọi thứ được sáng tỏ.''

Ryuuhou cũng quay sang phía chiếc đồng hồ, đồng lòng với chị. ``Mở hết toàn bộ tin nhắn của ngày hôm thứ Hai đó ra nhé, để chúng ta xem ai đúng ai sai!''

``Tôi sẽ chỉ truy xuất những dữ liệu thực sự cần thiết để giải quyết vấn đề này, không truy xuất thông tin thừa không liên quan,'' Game thông báo trước khi bắt đầu truy xuất dữ liệu.

``Mở đi,'' cả hai người đồng thanh nói lên, cùng nhau yêu cầu Game làm rõ sự việc.

Tôi nhìn sang Naomi, cô bé nhỏ bé cũng đang mở to mắt nhìn chằm chằm vào chiếc đồng hồ trên bàn, giống như mọi người khác trong phòng đều đang chờ đợi sự thật được phơi bày.

Game bật màn hình đồng hồ sáng bừng lên, chiếu toàn bộ lịch sử tin nhắn trong nhóm gia đình lên không trung nhỏ trước mặt hai người, từng dòng chữ hiện rõ ràng như đang mở một phiên tòa xét xử thật sự.

Tin nhắn cũ của chị Kaigou gửi từ hôm thứ Hai hiện lên đầu tiên, nội dung ngắn gọn: ``Ryuu, rửa giúp chị mấy cái nồi nhé, chị phải gọi cho mẹ.''

Ngay dưới đó là tin nhắn trả lời của Ryuuhou, chỉ vỏn vẹn bốn chữ: ``Vâng, em sẽ xử lý.''

Chị Kaigou lập tức khoanh tay trước ngực, giọng có chút thắng thế: ``Thấy chưa? Chị chỉ nhờ rửa mấy cái nồi thôi mà.''

Ryuuhou nhanh chóng chỉ ngón tay vào dòng chữ tin nhắn của mình, không chịu nhường bước: ``Em đã nói là sẽ `xử lý' mà. Em hiểu là phải lo hết cả phần bếp, chứ không chỉ có mỗi mấy cái nồi không đâu!''

``Em có thể hỏi lại chị nếu không rõ chứ!'' giọng Kaigou vẫn gắt.

``Thì chị cũng có thể viết rõ ràng hơn trong tin nhắn mà!'' Ryuuhou đáp trả ngay.

Game đột nhiên chen ngang vào cuộc cãi vã, giọng điện tử đều đều phá vỡ bầu không khí căng thẳng: ``Cụm từ `xử lý' không xác định rõ phạm vi công việc. Cả hai cách hiểu của Kaigou-user và Ryuuhou-user đều có cơ sở, không có cách nào kết luận ai đúng ai sai.''

Chị Kaigou lập tức quay sang nhìn chiếc đồng hồ, giọng đầy bực bội: ``Ông thật sự chẳng giúp được gì cả.''

``Tôi vừa xác nhận rằng yêu cầu của cô ban đầu thiếu độ chính xác cần thiết,'' Game giải thích một cách khách quan.

``Đấy, chị nghe Game nói chưa?'' Ryuuhou liền tận dụng cơ hội, giọng có chút vui mừng.

``Nhưng cho dù thế nào đi nữa, em vẫn chưa rửa cái bát của mình hôm nay.'' chị tôi lập tức quay về vấn đề gốc rễ của cuộc tranh luận này.

``Nhưng em đã rửa hết mấy cái nồi rồi mà!''

``Vậy thì cái bát vẫn còn nằm nguyên đó, em chưa động vào nó.''

Ryuuhou đột nhiên kéo ghế ra mạnh một chút, đứng dậy khỏi chỗ ngồi, cả người tỏ ra quyết tâm: ``Được thôi, em đi lấy thứ này để chị xem rõ.''

Nói rồi em ấy nhanh chóng rời khỏi phòng khách, bước chân vội vàng leo lên cầu thang gỗ dẫn lên các tầng trên.

Chị Kaigou đứng yên nhìn theo bóng dáng Ryuuhou biến mất ở cầu thang, rồi quay sang hỏi với vẻ khó hiểu: ``Em ấy định đi lấy gì thế?''

Tôi nhẹ nhàng đáp lời, vẫn giữ miếng vải lau trong tay: ``Có lẽ là lên phòng lấy gì đó để chứng minh thôi chị ạ.''

Naomi đang xếp khăn ăn bên cạnh cũng nghiêng đầu nhỏ của mình ra, đôi mắt tròn long lanh nhìn về phía cầu thang, em nhỏ giọng nói: ``Ryuuhou-oneechan có vẻ hơi giận thật sự.''

Chị Kaigou thở một hơi dài đầy mệt mỏi, vai cũng hơi xệ xuống ngán ngẩm: ``Chị đâu có muốn cãi nhau to tiếng thế. Chị chỉ muốn cái bát đĩa bẩn đó được rửa sạch thôi.''

Tôi biết chị nói hoàn toàn thật lòng, không có chút giả dối nào. Nhưng đồng thời, tôi cũng thấu hiểu rõ ràng rằng vấn đề bây giờ đã không còn đơn giản là một cái bát chưa rửa nữa. Nó đã trở thành cuộc đối đầu về những quan điểm khác nhau, về những cảm xúc bị dồn nén bấy lâu.

\ct

Chỉ vài phút sau, Ryuuhou quay lại phòng khách với một bảng kẹp giấy dày cộp trong tay, bìa bảng cứng cáp có vẻ đã được chuẩn bị từ lâu. Em ấy đặt cái bảng kẹp đó lên bàn ăn, ngay cạnh cuốn sổ nhỏ màu xanh của chị Kaigou, hai cuốn sổ đối diện nhau như hai đội quân đối đầu trên chiến trường.

Trên trang đầu tiên của bảng kẹp, một dòng tiêu đề lớn được viết bằng bút dạ màu xanh lá cây, nét chữ to và rõ ràng: ``Những việc em đã làm thêm trong nhà.''

Chị Kaigou nhíu mày nhìn cái bảng kẹp kỳ lạ, hỏi: ``Em cũng lập ra một cuốn sổ giống chị à?''

``Không phải sổ. Đây là bảng công trạng của em,'' Ryuuhou đáp một cách đầy tự hào.

``Vậy thì nó khác gì với cuốn sổ của chị đâu?''

``Cái tên của nó nghe đỡ giống sổ kỷ luật của trường học hơn chứ!''

Nói rồi em ấy lật sang trang tiếp theo của bảng, nơi có một danh sách dài được viết cẩn thận: ``Thứ Hai tuần này: em đã rửa hết mấy cái nồi, gom toàn bộ rác tái chế và tự mình mang cả thùng giấy lớn xuống tầng một để bỏ.''

Chị Kaigou nhìn vào dòng chữ đó, suy nghĩ một lúc rồi đáp: ``Mang thùng giấy xuống lẽ ra là việc của em hôm đó mà.''

``Nhưng trong lịch phân công công việc ghi đó là việc chung của hai người chúng ta mà,'' Ryuuhou nhanh chóng phản bác.

``Vì lúc đó chị tưởng hai đứa sẽ cùng nhau làm việc đó, chia sẻ công sức,'' chị Kaigou giải thích.

``Nhưng cuối cùng chỉ có mình em làm hết tất cả, một mình ôm cả cái thùng giấy nặng trịch xuống tầng một!'' Ryuuhou nói, giọng có chút oan ức.

Em ấy lại tiếp tục lật sang trang tiếp theo, chỉ vào dòng chữ tiếp theo: ``Thứ Ba: em đã gấp hết tất cả quần áo khô trong phòng khách vì máy sấy đột nhiên báo lỗi. Sau đó em cũng mang tất cả khăn sạch lên tầng sáu để mọi người dùng!''

Chị Kaigou lại khoanh tay, phản hồi lại: ``Máy sấy đã chạy lại bình thường sau khi em khởi động lại nó lần nữa mà nhỉ?''

``Vậy thì em vẫn phải bỏ thời gian ra để gấp hết đống quần áo đó chứ, công sức đâu phải tự nhiên mà có!''

``Đúng là vậy. Chị không bao giờ nói là em không làm những việc đó.''

``Nhưng chị cũng không bao giờ ghi lại những việc em đã làm vào cuốn sổ của chị. Tất cả những gì chị ghi là mỗi lần em quên làm gì đó.''

``Chị chỉ ghi những việc mà chị phải nhắc nhở em thôi, chứ không phải ghi tất cả mọi thứ trong nhà!'' 

``Chị chỉ ghi lại mỗi lần em chưa làm việc, còn những lần em làm thêm việc thì chị hoàn toàn không thấy!'' Ryuuhou nhẹ nhàng gõ ngón tay lên cái bảng kẹp, nhấn mạnh vào câu nói của mình.

Chị Kaigou đột nhiên im lặng hoàn toàn, không còn nói thêm lời nào.

Tôi nhìn chằm chằm vào cái bảng công trạng của Ryuuhou, lòng đầy ngạc nhiên. Danh sách những việc em ấy đã làm dài hơn rất nhiều so với những gì tôi tưởng tượng. Có một vài việc trong đó tôi biết rõ em đã làm, nhưng còn rất nhiều việc khác mà tôi hoàn toàn không hề hay biết.

Tôi nhận ra mình cũng từng nghĩ Ryuuhou chỉ làm những việc có tên trong lịch. Có lẽ em đã âm thầm nhận thêm nhiều việc hơn tôi thấy. Tôi khẽ nói: ``Ryuuhou-chan, chị không biết em làm thêm nhiều việc như vậy. Cảm ơn em nhé.''

\ct

Đúng lúc đó, tiếng chìa khóa xoay ổ cửa vang lên, cửa chính căn nhà bật mở ra một cách từ từ.

Kson và Delutaya bước vào cùng một lúc, cả hai người đều cầm trên tay một túi bắp rang bơ thơm phức, vừa mua từ cửa hàng tiện lợi dưới lầu. Họ vừa bước qua ngưỡng cửa thì mắt liền dính chặt vào khung cảnh kỳ lạ trong phòng bếp. Hai chị ấy lần lượt nhìn chiếc bồn rửa trong bếp còn ngổn ngang vài chiếc đĩa dơ, rồi nhìn cuốn sổ màu xanh nằm chình ình trên bàn, cuối cùng ánh mắt dừng lại ở bảng kẹp giấy dày cộp mà tôi vừa đặt xuống.

Kson thong thả bước đến bàn, đặt túi bắp rang của mình lên mặt gỗ còn hơi ẩm vệt nước lau. Giọng chị ấy tràn đầy sự ngạc nhiên, như thể vừa bước nhầm vào một buổi họp cổ đông bất ngờ: ``Bọn chị vào nhầm thời điểm à?''

Delutaya cũng hoàn tất việc tháo giày, xếp gọn đôi sneaker vào kệ giày sát tường ở tầng dưới, rồi quay lại liếc qua hai chị em Kaigou và tôi, mỗi người đang đứng một đầu bàn với tài liệu của mình. Chị ấy cũng gật gù đồng tình: ``Có vẻ như đúng là đang có một cuộc họp gì đó căng thẳng lắm. Mọi người trông nghiêm túc quá.''

Tôi thở dài đáp lại, nhíu mày: ``Hai người họ lại tranh cãi nhau vì một cái bát ạ.''

``Thật á?''

``Dạ.''

Bên kia bàn, chị Kaigou nhanh chóng giơ cuốn sổ màu xanh của chị ấy lên, không chịu kém cạnh: ``Còn đây là lịch trực nhật việc nhà chính thống, ghi rõ ràng từng đầu việc ai nhận, ai hoàn thành. Không phải cứ liệt kê bừa bãi như ai đó.''

Kson vừa mở miệng túi bắp rang, tiếng xốp giòn vang lên trong phòng yên tĩnh. Chị nhặt một hạt bắp bỏ vào miệng, nhai vài cái rồi nói với giọng trầm trồ: ``Hai cái tên nghe hay nhỉ, gần nhau thật đấy. Công bố thành tích với lịch trực nhật, nghe đâu như hai bộ phận cùng làm việc trong một công ty lớn.''

``Không hề giống nhau chút nào,'' chị Kaigou đáp ngay, giọng vẫn còn gắt gỏng. ``Cái của tôi là để quản lý công việc, còn cái bảng kia thì\dots''

``Không đâu!'' Ryuuhou ngắt lời chị, vẫn giữ nguyên vẻ mặt cứng đầu. ``Thành tích của em ghi rõ từng việc tôi đã làm thêm ngoài lịch, còn chị thì chỉ biết ghi mỗi lỗi của người khác! Hai cái hoàn toàn khác bản chất.''

Delutaya nhìn hai túi bắp rang còn nguyên trên bàn, rồi lại nhìn bộ dạng căng thẳng của hai chị em, ánh mắt tràn đầy sự bối rối. Chị nói nhỏ: ``\hl{Hai người họ trông có vẻ bướng bỉnh ghê\dots}''

Kson cũng nhìn qua chiếc bồn rửa trong bếp, rồi nhìn lại hai chúng tôi, cười nhạo một cách tinh nghịch: ``\hl{Nhưng chị thấy cái này còn gay cấn hơn cả phim hành động mà ta định xem nữa.}''

Tôi nghe vậy cũng chỉ có thể cười gượng với lời bình của ``Chị Đại ngổ ngáo'' ấy.

Kaigou liền chỉ tay vào túi bắp rang trên bàn, giọng đầy cảnh cáo: ``Đừng có ngồi đó xem trò vui nhé! Chuyện này thật sự nghiêm túc đấy!''

Kson lập tức đặt tay lên túi bắp rang, vẻ mặt vô tội như vừa bị bắt quả tang làm gì sai: ``Em không định xem đâu, thật lòng đấy. Em chỉ muốn hỏi xem hai bên có cần trọng tài để phân xử đúng sai không? Tránh việc cãi nhau mãi không ra kết quả.''

Ngay lúc đó, chiếc đồng hồ thông minh trên bàn bỗng phát ra tiếng bíp, màn hình sáng bừng lên. Game - người bạn đồng hành kỹ thuật số của chúng ta - lên tiếng với giọng điện tử đều đều: ``Tôi có thể đảm nhiệm vai trò trọng tài hoàn toàn khách quan. Dữ liệu của tôi có thể làm rõ mọi sự việc, không thiên vị bất kỳ bên nào.''

Kaigou và Ryuuhou lập tức đồng thời quay sang nhìn chiếc đồng hồ, cùng nhau mở miệng nói một câu: ``Không cần!''

Delutaya thấy vậy liền kéo ghế trong phòng khách ra, ngồi xuống một cách thoải mái, đặt túi bắp rang của mình lên đùi. Chị ấy nhún vai, chấp nhận số phận phải ngồi đây: ``Vậy thì bọn chị chỉ ngồi đây để cổ vũ tinh thần thôi. Không can thiệp vào cuộc tranh luận của các em. Chỉ đứng ngoài xem thôi, thật lòng.''

Tôi nhìn túi bắp rang còn nguyên trên bàn, rồi nhìn bộ dạng của tất cả mọi người trong phòng: chị Kaigou vẫn đang cau có với cuốn sổ của mình, Kson đang chăm chú nhìn chằm chằm vào hai chúng tôi như xem chương trình trực tiếp, Delutaya cũng đã mở túi bắp rang lấy một ít bỏ vào miệng, Naomi ở góc phòng cũng đang mở to mắt nhìn về phía bàn. 

Lúc đó tôi thầm nghĩ trong lòng: buổi tối này chắc chắn sẽ còn kéo dài rất dài, không thể kết thúc sớm được.

\ct

Cuộc tranh cãi giữa hai chị em bùng nổ đầy kịch tính ngay giữa căn bếp, buộc tôi cùng với Kson và Delutaya - những người chỉ tình cờ có mặt vào đúng lúc này - phải đóng vai những khán giả vô tình, không thể nào rời đi mà phải chứng kiến toàn bộ cuộc đối đầu căng thẳng giữa Kaigou và Ryuuhou.

Chị Kaigou chính thức mở đầu cho phần phản bác của mình, giọng điệu vẫn giữ được sự bình tĩnh nhưng đã có những dấu hiệu của sự căng thẳng khó che giấu. Chị cầm cuốn sổ màu xanh lam đã góc bìa hơi mòn, lật từng trang nhẹ nhàng cho đến khi dừng lại ở một tờ giấy vàng nhạt đánh dấu trang, đó là những trang chị thường xuyên mở ra nhất.

``Chị không bao giờ ghi lại mọi lần em quên làm việc. Những gì chị ghi vào sổ chỉ là những lần em đã nhận trách nhiệm với một công việc cụ thể, nhưng lại bỏ dở giữa chừng.''

Ryuuhou vừa nghe vừa cúi đầu nhìn vào bảng kẹp giấy của mình, những dòng chữ viết tay của em trải rộng trên trang giấy trắng, từng mục việc làm thêm được ghi cẩn thận: ``Em cũng không phải ghi lại mọi thứ em đã làm. Em chỉ liệt kê những công việc em đã tình nguyện nhận thêm, những việc không nằm trong lịch phân công ban đầu, mà không ai ghi nhận hay tính vào trách nhiệm của em.''

% ``Mà em cũng không hề nói với chị là em đã nhận thêm những việc đó.'' Kaigou đáp lại ngay, vẫn giữ nguyên giọng điệu.

% ``Lúc đó chị đang ở ngoài ban công nói chuyện điện thoại với mẹ, em có gọi mà chị không nghe thấy.'' Ryuuhou giải thích, ngón tay vô tình vo nhẹ góc bảng kẹp.

% ``Thì em có thể nhắn tin cho chị mà.''

% ``Em tưởng chị sẽ tự nhìn thấy em đang làm, nên không nghĩ phải nhắn tin riêng.''

% ``Chị làm sao có thể nhìn thấy tất cả mọi việc em làm nếu em không tự báo cho chị biết? Chị cũng có công việc của mình phải lo, không thể lúc nào cũng để ý xem ai đang làm gì trong nhà.''

% ``Thì em cũng làm sao biết được chị vẫn tính rằng phần việc đó là của em, nếu chị không tự nói với em? Chúng ta cứ thế hiểu lầm qua nhau, làm sao mà rõ ràng được.''

Ở phía ghế sofa, Kson khẽ thổi một tiếng huýt sáo nhỏ, không muốn làm phiền cuộc đối thoại nhưng cũng không giấu được sự ngạc nhiên trước cách hai người đẩy đi đẩy lại vấn đề.

Delutaya cũng vươn tay lấy một nắm bắp rang bơ từ túi trên đùi, nhưng hạt bắp vẫn nằm yên trong lòng bàn tay, chị ấy chẳng thể nào bỏ vào miệng được khi bầu không khí trong phòng đang căng thẳng đến thế.

Ryuuhou đột nhiên đưa ngón tay chỉ vào một dòng khác trong bảng công trạng của mình, đó là một việc được em gạch chân rõ ràng.

``Thứ Tư tuần này, em đã tự mình mang cả thùng rác tái chế lớn xuống tầng trệt để bỏ. Lịch phân công vốn ghi tên em phải làm việc đó, nhưng em còn làm luôn cả phần của chị nữa, vì hôm đó chị đang bận hoàn thành báo cáo công việc ở nhà.''

Kaigou trả lời ngay lập tức, không chần chừ gì: ``Chị đã xin em giúp chị làm phần đó mà.''

``Nhưng lúc đó chị nói là `để chị tự làm sau, em đừng lo', chị đâu có bảo em làm thay chị đâu?''

``Rồi cuối cùng chị cũng đã tự làm những việc còn lại mà, đâu phải em làm hết tất cả.''

``Nhưng em đã làm trước chị rồi, em thấy thùng rác đầy quá nên mang đi trước, có sao đâu?''

``Vậy thì em đã giúp chị hoàn thành việc đó, chị rất biết ơn.''

``Nhưng chị đâu có bao giờ tính đó là phần việc em đã làm thêm, chị chỉ nhớ mỗi lần em quên làm phần việc của mình thôi.''

Kaigou khẽ cau mày, hai lông mày chị xích lại gần nhau, vẻ mặt có chút bối rối trước lời nói của em gái: ``Chị thật sự không biết là em muốn được tính công cho những việc đó.''

``Em không muốn được tính công gì cả, em chỉ không muốn bao giờ phải nghe chị nói rằng em chẳng bao giờ làm gì trong nhà, như thể em chỉ biết ăn không ngồi rồi!''

Lời nói của Ryuuhou khiến Kaigou sững lại hoàn toàn, chị đứng yên ở đó, không nói thêm được câu nào. Chị chưa bao giờ nói ra câu đó, chưa bao giờ nghĩ rằng em gái sẽ hiểu lầm chị đến thế.

Nhưng tôi đứng yên ở góc bếp, vẫn giữ nguyên miếng vải lau bàn trong tay, mắt theo dõi từng biến đổi nhỏ trên khuôn mặt của hai chị em và lắng nghe từng lời đối đáp của họ. Tất cả những gì diễn ra trước mắt tôi, từ sự bực bội của Kaigou cho đến nỗi oan ức dồn nén bấy lâu của Ryuuhou, đều làm tôi thấu hiểu hoàn toàn những gì em gái đang cố gắng truyền tải. 

Tôi cảm nhận rõ ràng sự tổn thương trong giọng nói của Ryuuhou, cảm giác như những công sức em bỏ ra mỗi ngày để giữ cho căn nhà ngăn nắp đều bị bỏ qua, không ai ghi nhận hay thậm chí nhận ra sự tồn tại của chúng. Cái cảm giác bị đánh giá chỉ qua những lỗi lầm nhỏ, bị nhìn nhận như một người luôn trốn tránh trách nhiệm, mà không ai nhớ đến những lần em tự giác đứng dậy làm thêm việc, đã ăn sâu vào lòng em như một vết xước dai dẳng.

Ryuuhou đột nhiên kéo ghế ra, ngồi xuống đối diện với chị Kaigou, hai người đối mặt nhau trên bàn ăn: ``Vậy thì tại sao chị lại phải giận dữ đến thế chỉ vì một cái bát chưa rửa?''

Kaigou mở miệng định nói gì đó, rồi lại khép lại, chị không tìm được lời nào để đáp lại ngay lập tức. Chị nhìn xuống cuốn sổ màu xanh lam trong tay mình, những dòng chữ viết tay của chị hiện rõ trước mắt.

``Vì chị cũng mệt.''

Giọng nói của Kaigou đột nhiên thấp hẳn đi, không còn chút gắt gỏng nào nữa, chỉ còn lại sự mệt mỏi dồn nén bấy lâu.

``Chị không phải bực mình chỉ vì một cái bát bẩn. Chị bực vì cứ mãi nghĩ rằng nếu chị không nhắc nhở từng người một, thì tất cả mọi việc sẽ bị bỏ quên ở đó, không ai động vào. Chị ghi chép vào sổ vì sợ mình quên mất ai đã nhận việc gì, nhưng cuối cùng mỗi lần mở sổ ra nhìn những dòng chữ, chị lại cảm thấy như mọi thứ chẳng bao giờ thay đổi, cứ lặp đi lặp lại mãi.''

Ryuuhou không đáp lại ngay, em cũng đứng im lặng, xử lý những gì chị vừa nói.

Kaigou nhìn em gái, giọng càng lúc càng nhẹ: ``Nhưng chị thật sự không biết em đã làm tất cả những việc kia, nếu em không nói ra thì chị làm sao mà biết.''

``Đúng là em cũng chưa bao giờ nói với chị về những việc đó.'' Ryuuhou thừa nhận.

``Chị hiểu rồi.''

``Em cứ nghĩ rằng nếu em tự giác làm việc đó, thì chị sẽ tự nhiên nhìn thấy, nên không cần phải nói ra để khoe công.''

Kaigou thở một hơi dài, như thể vừa thả hết được tất cả sự căng thẳng trong người từ nãy giờ.

``Chị đâu có phải là Game, chị không thể theo dõi từng giây từng phút mọi việc xảy ra trong nhà, không thể thấy hết tất cả những gì mọi người làm.''

Đúng lúc đó, chiếc đồng hồ thông minh trên bàn lại lên tiếng, giọng điện tử đều đều phá vỡ khoảnh khắc bình yên vừa xuất hiện: ``Tôi cũng không theo dõi được mọi việc trong nhà. Có một số khu vực trong nhà không nằm trong vùng phủ sóng cảm biến, tôi cũng không thể ghi nhận được.''

Ryuuhou lập tức quay sang nhìn chiếc đồng hồ, giọng có chút bất lực: ``Ông đừng có phá hỏng khoảnh khắc vừa rồi của chúng tôi được không?''

``Tôi chỉ muốn làm rõ giới hạn về dữ liệu mà tôi có thể thu thập được thôi, để mọi người không hiểu lầm thôi mà\dots'' Game giải thích.

\ct

Chị Kson đã nhai gần như cạn sạch túi bắp rang bơ đầu tiên, chỉ còn lại vài hạt lẻ loi nằm ở góc túi. Còn Delutaya thì túi thứ hai vẫn còn nguyên vẹn, chưa hề động tay vào một hạt nào; ánh mắt của chị quá tập trung vào sự đối đầu căng thẳng trước mặt mà quên mất món đồ ăn đang đợi trên đùi.

Hai người họ ngồi chễm chệ trên ghế gỗ, đôi mắt không ngừng đảo qua đảo lại liên tục giữa cuốn sổ màu xanh lam cũ mòn của Kaigou, và bảng kẹp giấy dày cộp với dòng tiêu đề ``bảng công trạng'' của Ryuuhou đầy đối lập nhau ở trong một phiên tòa.

``Em thấy cả hai bên đều có đủ bằng chứng để bảo vệ quan điểm của mình,'' Kson lên tiếng phá vỡ im lặng, giọng trầm trồ như một nhà bình luận đang theo dõi một trận đấu bóng đá kịch tính.

Kaigou lập tức đáp trả ngay, giọng vẫn còn gắt gỏng: ``Đây không phải tòa án để mà mang bằng chứng ra trình chiếu đâu nhé.''

``Thế mà sao mọi người trong phòng này ai cũng cầm một xấp tài liệu dày cộp thế kia? Không giống tòa án thì giống gì?'' chị tóc tím nhướng mày phản vấn lại, cầm thêm hạt bắp rang ném vào miệng.

Ryuuhou nhanh chóng chìa chiếc bảng kẹp giấy của mình ra, phản bác lại lời nói của Kson: ``Em chỉ có mỗi một bảng kẹp giấy này thôi mà! Không phải tài liệu gì đâu.''

``Đúng đúng, bằng chứng có kẹp cả vào bảng thế kia mà,'' Kson liền châm biếm, lời chơi chữ làm không khí vốn đã căng thẳng càng thêm phần dở khóc dở cười.

Ngay bên cạnh, Delutaya vội vàng kéo nhẹ nhàng vào tay áo Kson, giọng nói nhỏ như sợ ai nghe thấy: ``Chị đừng châm thêm dầu vào lửa nữa được không ạ? Mọi thứ đang căng thẳng lắm rồi.''

Kson đành phải đặt nắm bắp rang đang cầm trong tay xuống ghế, giơ hai tay lên ra vẻ đầu hàng: ``Được thôi, chị chỉ đang cố gắng giữ cho bầu không khí nó được nhẹ nhàng một chút thôi mà. Làm gì mà căng thế.''

``Nhưng bầu không khí giờ đang nặng hơn cả cái thùng giấy tái chế em mang xuống tuần trước rồi chị ơi,'' chị tóc xanh lắc đầu ngao ngán, mắt vẫn không rời khỏi hai chị em đang đối mặt nhau ở đầu bàn.

Đúng lúc đó, chiếc đồng hồ thông minh trên bàn đột nhiên bật màn hình sáng lên to hơn bao giờ hết, phá vỡ cuộc đối thoại nhỏ của hai người trên ghế sofa. Giọng điện tử đều đều của Game vang lên, với vẻ điệu đà như một chuyên gia dữ liệu sắp trình bày báo cáo tài chính: ``Tôi đã hoàn thành việc tổng hợp toàn bộ lịch phân công việc nhà của cả gia đình trong bảy ngày gần nhất. Dữ liệu đã được kiểm tra và xác minh đầy đủ.''

Kaigou lập tức quay người sang phía chiếc đồng hồ, không chần chừ gì nữa mà ra lệnh ngay: ``Chiếu hết ra đi. Để mọi thứ được minh bạch rõ ràng.''

Ryuuhou cũng gật đầu liên tục, đồng lòng với chị gái: ``Chiếu luôn đi, để ai đúng ai sai cũng được sáng tỏ.''

Game liền hiển thị lên màn hình lớn một biểu đồ cột rõ ràng, được chia thành hai phần chính nhau. Một cột màu xanh dương ghi lại tất cả số công việc theo lịch phân công đã được hoàn thành đúng hạn. Còn cột còn lại màu cam ghi lại số công việc đã được làm thêm ngoài lịch phân công ban đầu, những việc mà không ai giao mà tự nguyện làm.

Kaigou lập tức chăm chú nhìn vào cột màu xanh dương ở đầu tiên, còn Ryuuhou thì mắt dính chặt vào cột màu cam thứ hai, ai cũng tìm ra điểm mạnh của mình trên bảng dữ liệu.

``Tổng số lượt rửa bát của Kaigou-user trong bảy ngày là bốn lượt. Ryuuhou-user là hai lượt. Về các công việc ngoài lịch, Ryuuhou-user có ba lượt tình nguyện làm thêm, còn Kaigou-user có hai lượt,'' ông ấy đều đều đọc ra những con số trên biểu đồ, không có chút cảm xúc nào.

Ryuuhou chớp mắt vài lần, như thể không tin vào những con số mình vừa nghe thấy. Liền chỉ ngón tay vào cột màu cam trên màn hình, em quay sang nói với Kaigou, giọng có chút phấn khích: ``Thấy chưa? Em làm việc ngoài lịch nhiều hơn hẳn chị ấy! Em đã làm thêm nhiều việc mà.''

Kaigou không chịu thua, chỉ mạnh vào cột màu xanh dương trước mặt mình, đáp lại ngay: ``Nhưng chị hoàn thành tất cả những phần việc được giao cho mình nhiều hơn em. Còn em thì để sót cả cái bát hôm nay trong bồn rửa nữa.''

Game lại tiếp tục đọc thêm những dữ liệu khác, làm tăng thêm sự phức tạp của vấn đề: ``Trong bảy ngày qua, có tổng cộng ba lần thay đổi phân công việc nhà được nhắc đến trong nhóm tin nhắn của gia đình. Trong đó có hai lần thay đổi không nhận được xác nhận từ người nhận việc, và một lần có xác nhận nhưng không được cập nhật vào lịch chung của cả nhà.''

Kson từ trên ghế sofa ngó sang nhìn biểu đồ, suy nghĩ một lúc rồi đưa ra nhận xét của mình: ``Vậy kết luận là hai người đều đúng một nửa. Không có ai hoàn toàn sai, cũng không có ai hoàn toàn đúng.''

Kaigou lập tức lắc đầu liên tục, không đồng tình với nhận xét đó: ``Không hề. Vấn đề cốt lõi vẫn là em ấy để cái bát bẩn nằm nguyên trong bồn rửa mà không chịu rửa. Đó là sự thật.''

Ryuuhou cũng chỉ mạnh vào cột dữ liệu của mình trên màn hình, đáp trả lại chị gái: ``Còn vấn đề của em là chị không bao giờ ghi nhận hay tính đến những việc em đã làm thêm ngoài lịch đó! Chỉ nhớ mỗi lần em quên làm phần việc của mình thôi!''

Game đột nhiên phóng to hai dòng dữ liệu nhỏ ở cuối biểu đồ, giọng vẫn đều đều không đổi: ``Tôi không thể xác định rõ ý nghĩa chính xác của cụm từ `đổi ca' mà hai người đã nhắc đến trong tin nhắn, nếu hai người không ghi rõ ràng thời hạn và phạm vi cụ thể của công việc được đổi. Dữ liệu này không đủ để đưa ra kết luận.''

Kaigou nhìn chằm chằm vào màn hình đồng hồ, giọng đầy bực bội: ``Ông nói chuyện như thể tất cả lỗi lầm này đều là của cái biểu mẫu lịch phân công vậy!''

``Biểu mẫu lịch phân công hiện tại của chúng ta thực sự có những thiếu sót nhất định, không đáp ứng được nhu cầu ghi nhận tất cả các thay đổi phát sinh được,'' ông ấy trả lời một cách khách quan, đúng như những gì dữ liệu cho thấy.

Ryuuhou đứng ở bên kia bàn, nhíu mày suy nghĩ những gì Game vừa nói. Liền em lên tiếng, giọng nói nhỏ hơn trước, như thể vừa tỉnh ngộ ra một điều gì đó: ``Nhưng cuối cùng thì bọn em mới là người phải nói chuyện rõ ràng với nhau. Dù biểu mẫu có tốt đến đâu mà chúng ta không trao đổi thì cũng vô dụng thôi.''

Tôi nhìn chiếc biểu đồ thu nhỏ trên màn hình. Game đã tìm được những chỗ thiếu rõ ràng trong tin nhắn, nhưng không thể thay hai chị em nói ra điều họ thật sự cần. Tôi khẽ lên tiếng: ``Em nghĩ Ryuuhou-chan nói đúng. Mình vẫn cần nghe chính hai người giải thích.''

Ngay sau câu nói đó, không có ai trong phòng đáp lại lời em. Cả Kson và Delutaya trên ghế sofa cũng im lặng, Naomi ở góc phòng cũng không phát ra tiếng động nào.

Và đây là lần đầu tiên kể từ khi cuộc cãi nhau bắt đầu, cả hai chị em Kaigou và Ryuuhou cùng đồng thời cúi nhìn xuống những ghi chép của mình trên bàn, nhìn vào cuốn sổ màu xanh và bảng kẹp giấy đã làm họ đối đầu nhau suốt buổi tối.

Sự im lặng lần này khác hẳn lúc đầu. Không còn là khoảng dừng trước câu đáp trả kế tiếp, mà giống như cả hai đang lần đầu nhìn vào những gì đối phương đã cố ghi lại.

\ct

Đúng lúc mọi người đang dần bình tĩnh lại, tiếng cánh cửa trượt nặng nề từ từ mở ra, và anh Kuon bước vào, gót giày chạm nhẹ vào sàn gỗ của căn nhà. Ánh mắt anh quét qua toàn bộ khung cảnh trước mặt, và có lẽ trong giây lát, anh cũng ngạc nhiên không kém gì Kson và Delutaya khi mới bước vào.

Ở phía đối diện, Kaigou vẫn đang đứng nghiêm sau bàn ăn, cuốn sổ màu xanh lam cũ mòn của chị mở rộng ra trước mặt, những dòng chữ viết tay rõ ràng hiện ra. Ngồi đối diện chị, Ryuuhou cũng không nhúc nhích, chiếc bảng kẹp giấy dày cộp với dòng tiêu đề "bảng công trạng" vẫn trải phẳng trên mặt bàn, như một bản tuyên ngôn không lời. Trên chiếc đồng hồ thông minh ở giữa bàn, Game vẫn đang chiếu cái biểu đồ cột rực rỡ, những con số hiện lên sáng lấp lánh trong ánh đèn vàng ấm của phòng khách. 

Trên chiếc sofa lớn ở góc phòng, Kson và Delutaya vẫn ngồi đó, túi bắp rang bơ đã mở nhưng chỉ còn vài hạt lẻ loi, hai chị ấy vẫn đang chăm chú theo dõi như hai khán giả trung thành của một buổi biểu diễn. Còn Naomi, bé gái nhỏ nhất trong nhà, đang ngồi xếp khăn ăn trên một góc bàn, xếp thành một hàng thẳng tắp, ngay ngắn không khác gì một hàng lính nhỏ đang chờ duyệt binh.

Không ai nói gì trong vài giây im lặng đó, cho đến khi Kson phá vỡ sự tĩnh lặng, giọng chị vẫn mang chút tinh nghịch như mọi khi: ``Kuon-paisen về đúng giờ thật đấy, vừa kịp lúc phần cao trào nhất.''

Kuon nhướng mày, đặt chìa khóa vào túi quần rồi hỏi: ``Có ai gọi anh về sớm để giải quyết chuyện gì không? Tại sao bầu không khí căng thẳng thế này?''

Cùng một lúc, Kaigou và Ryuuhou cùng mở miệng, cùng nhau nói ra một câu: ``Không có ai gọi anh cả.''

Chiếc đồng hồ thông minh cũng thêm vào, giọng điện tử đều đều không thay đổi: ``Tôi cũng đã không gửi thông báo nào đến Kuon-user. Không có yêu cầu hỗ trợ khẩn cấp nào được ghi nhận.''

Kuon bước lại gần bàn, nhìn cái biểu đồ đang chiếu trên màn hình đồng hồ, những cột xanh dương và cam nhô lên rõ rệt. Anh thở một hơi dài, đầy ngạc nhiên: ``Vậy sao mà anh có cảm giác mình vừa bước chân vào một cuộc họp cổ đông khẩn cấp của một công ty sắp phá sản thế này? Mọi người ai cũng cầm tài liệu, ai cũng trông nghiêm túc quá.''

``Bọn em đang giải quyết chuyện phân công việc nhà trong gia đình thôi ạ,'' Kaigou đáp ngay, giọng chị vẫn còn chút gắt gỏng chưa tan hết.

``Giải quyết bằng cách lập hồ sơ hai phía rồi trình bằng chứng ra trước tòa à?'' anh Kuon cười nhạt, chỉ tay vào cuốn sổ và cái bảng kẹp đang đối đầu nhau trên bàn.

Ryuuhou liền kéo chiếc bảng công trạng của mình về phía mình, như thể bảo vệ tài liệu quý giá, rồi nói: ``Đây không phải hồ sơ gì cả! Đây là danh sách những việc em đã làm thêm ngoài lịch phân công ban đầu đó, Onii!''

``Còn đây là lịch trực nhật chính thức của cả nhà,'' Kaigou cũng nhanh chóng bổ sung, giơ nhẹ cuốn sổ màu xanh lên một chút.

Kuon không nói gì thêm, chỉ đứng yên ở cửa bếp, nhìn vào chiếc bát duy nhất còn nằm lại trong bồn rửa, nguyên nhân gốc rễ của toàn bộ cuộc tranh cãi đêm nay. Rồi anh nhìn lại bề mặt bàn ăn, nơi có vài vệt nước còn sót lại sau khi lau, và một chồng đĩa sạch chưa kịp xếp vào giá.

Sau khi ngắm nghía đủ, anh quay ra hỏi hai chị em: ``Hai người đã thống kê xong tất cả số liệu chưa? Đã có kết luận cuối cùng chưa?''

Kaigou đáp ngay, chỉ vào chiếc đồng hồ: ``Game vừa mới tổng hợp xong tất cả số liệu, vừa bây giờ thôi.''

``Anh không hỏi Game,'' anh lắc đầu, giọng anh bình thản nhưng đầy thuyết phục. Anh quay lại nhìn thẳng vào mặt hai chị em, ánh mắt chân thành: ``Anh hỏi hai người. Hai người đã nói chuyện với nhau, nói ra hết những gì mình muốn nói chưa? Đã hiểu được ý của người kia chưa?''

Tôi thấy nhẹ nhõm khi Kuon-sama không chọn một cuốn sổ để làm trọng tài. Anh đang hỏi hai chị em xem họ đã thật sự nghe nhau chưa, chứ không hỏi ai có nhiều dòng ghi chép hơn.

Ryuuhou khoanh tay trước ngực, vẻ mặt vẫn còn chút cứng đầu, em đáp nhỏ: ``Chưa ạ.''

Kaigou cũng gật đầu đồng tình, chị cũng thừa nhận: ``Chúng em vẫn đang nói chuyện.''

Kuon gật đầu như đã đoán trước được câu trả lời. Anh kéo một chiếc ghế gỗ ra từ dưới bàn, ngồi xuống cạnh Naomi, cô bé vừa mới xếp xong hàng khăn ăn. Anh đặt hai tay lên đầu gối, nói một cách đơn giản: ``Vậy thì cứ tiếp tục nói đi. Anh sẽ ngồi đây nghe. Có bao nhiêu tâm can trút hết đi..''

Không ai trong phòng phản đối lời anh. Cả Kson và Delutaya trên sofa cũng im lặng, Naomi cũng không nhúc nhích. Và từ lúc đó, không ai còn nhìn vào chiếc đồng hồ thông minh nữa, không ai quan tâm đến những con số trên biểu đồ nữa. Tất cả ánh mắt đều quay lại với hai chị em, những người thực sự cần nói chuyện với nhau.

Chị Kaigou là người đầu tiên phá vỡ sự im lặng sau khi Kuon ngồi xuống. Chị thở dài, cơn bực bội ban đầu đã dần tan biến, thay vào đó là sự mệt mỏi hiện rõ trên mặt: ``Chị không bao giờ nghĩ em chẳng làm gì cả trong nhà. Chị chỉ bực mình vì những việc mà em đã nhận trách nhiệm thì lại không hoàn thành đúng hạn, và chị lại không hề biết những phần nào em đã đổi với người khác, không có thông báo gì cả.''

Ryuuhou cúi đầu nhìn những dòng chữ trên cuốn sổ màu xanh của chị, những dấu sao đỏ nhỏ mà Kaigou đã ghi lại bên cạnh những việc chưa hoàn thành. Em nói nhỏ, giọng không còn gắt gỏng như trước: ``Em không cố tình bỏ dở bất cứ việc nào đâu. Có những hôm em nghĩ là mình đã đổi ca với chị thành công, có hôm em định làm sau rồi tự dưng quên mất. Nhưng cũng có những lúc em đã tự giác làm thêm những việc ngoài lịch, những việc không phải của em, mà chị cũng chẳng bao giờ hay biết.''

Kaigou gật đầu chậm rãi, như thể đang thấu hiểu từng lời em nói. Chị nói một cách nhẹ nhàng: ``Vậy thì những gì hai đứa nói đều đúng hết. Cả hai đều có lý do của mình.''

Ryuuhou ngước lên, mắt em hơi ngạc nhiên, em không nghĩ chị sẽ thừa nhận điều đó. Em hỏi: ``Chị không định nói em sai hết, tất cả lỗi của em à?''

``Chị đang nói rằng chị không thể biết em đã làm những gì nếu em không tự nói với chị. Chị không phải là người có thể theo dõi mọi thứ mọi lúc.''

Kuon lúc này đặt hai khuỷu tay lên mặt bàn, nghiêng người về phía hai chị em, góp lời vào cuộc nói chuyện: ``Vậy vấn đề cốt lõi không phải là ai làm nhiều hơn ai, ai làm ít hơn ai. Cái đó không quan trọng bằng.''

Hai chị em cùng lúc quay sang nhìn anh, chờ đợi anh nói tiếp. 

``Câu hỏi thực sự là làm sao để người nhận việc biết rõ phần nào là trách nhiệm của mình, và khi ai đó muốn đổi ca, thì phải đảm bảo rằng người kia đã đồng ý rõ ràng, không có hiểu lầm nào xảy ra ả,'' anh nói đơn giản, nhưng chạm đúng vào gốc rễ của toàn bộ tranh cãi.

Ngay lúc đó, Game thu nhỏ cái biểu đồ trên màn hình lại, giọng điện tử của ông ấy lại vang lên: ``Tôi có thể tạo một biểu mẫu trực tuyến để xác nhận việc đổi ca. Tất cả các thay đổi đều cần có chữ ký số của cả hai bên để được cập nhật vào lịch chung. Điều này sẽ tránh được các hiểu lầm.''

Kuon nhìn chiếc đồng hồ, gật đầu nhưng cũng nhấn mạnh: ``Cái đó có thể làm sau. Nhưng trước hết, hai người phải tự thống nhất được cách làm việc với nhau, nói chuyện rõ ràng với nhau. Công nghệ không thể thay thế được việc hai người ngồi lại và nói chuyện được đâu.''

Kaigou lúc này gập cuốn sổ màu xanh của mình lại, đặt nó lên mặt bàn. Chị nhìn thẳng vào em gái, nói rõ ràng: ``Chị chỉ muốn rằng nếu em không thể hoàn thành phần việc của mình đúng giờ, thì hãy nói với chị sớm. Để chị biết mà sắp xếp lại, không phải cứ chờ đợi vô ích.''

Ryuuhou cũng gật đầu liên tục, em cũng nói ra mong muốn của mình: ``Còn em chỉ muốn rằng trước khi chị ghi em là bỏ dở việc, là em không hoàn thành trách nhiệm, thì hãy hỏi em một câu. Đừng vội ghi nhận ngay ạ.''

``Nếu em nhắn tin vào nhóm nói là em muốn đổi ca, thì chị sẽ trả lời xác nhận ngay lập tức. Thế thôi.''

``Nhưng nếu chị không trả lời tin nhắn của em thì sao? Thì việc đó vẫn là của ai?''

``Thì nếu không có xác nhận, phần việc đó vẫn thuộc về người ban đầu nhận. Không có gì thay đổi cho đến khi hai bên đều đồng ý.'' Kaigou đáp ngay.

``Vậy thì mọi thứ đều rõ ràng rồi,'' Ryuuhou mỉm cười nhẹ, cuối cùng cũng cảm thấy mọi thứ được sắp xếp ổn thỏa.

Kaigou nhìn xuống chiếc bảng công trạng của Ryuuhou, chị hỏi thêm: ``Còn những việc em làm thêm ngoài lịch phân công thì sao? Em muốn chị ghi chúng vào đâu không?''

``Em sẽ báo vào nhóm chung khi em làm thêm việc gì đó. Không cần phải ghi thành tích, không cần làm gì cả. Chỉ cần mọi người biết là em đã làm là được,'' Ryuuhou nói. Em ngập ngừng một lúc, rồi thêm vào, giọng thật nhỏ. ``Em không lập cái bảng công trạng này vì muốn được khen ngợi, muốn được ai đó tán dương. Em chỉ muốn chị biết rằng em cũng có đóng góp vào việc nhà, em không phải là người chỉ biết ăn không ngồi rồi.''

Tôi cúi nhìn chiếc khăn ăn mình vẫn đang xếp. Tôi cũng thường tự làm nốt những việc còn lại trong bếp mà không nói với ai, rồi nghĩ rằng như vậy sẽ đỡ phiền mọi người. Có lẽ sự im lặng của tôi đôi khi cũng khiến người khác khó biết mình đang cần giúp đỡ hay chỉ muốn tự lo. Tôi khẽ nói: ``Em cũng sẽ báo cho mọi người khi nhận thêm việc, để không ai phải đoán nhé.''

Kaigou dịu giọng hoàn toàn, chị nhìn em gái mình với ánh mắt thương hại: ``Chị biết rồi. Bây giờ chị đã hiểu hết những gì em muốn nói.''

\ct

Sau khi mọi thứ đã được nói rõ, Kuon đứng dậy từ ghế, bước lại vào gian bếp một lần nữa. Anh nhìn vào bồn rửa, nơi chiếc bát bẩn vẫn còn nằm đó, rồi nói vọng ra: ``Tốt lắm. Vậy là mọi việc đã được giải quyết. Nhưng phần việc trước mắt vẫn còn đó, hai người tự giải quyết với nhau nha.''

Kaigou quay sang nhìn Ryuuhou, chị nói: ``Vậy em rửa cái bát của em nhé, cái bát mà hôm nay em để lại trong bồn.''

Ryuuhou nhìn vào chồng bát đĩa sạch chờ xếp, rồi hỏi: ``Còn những cái còn lại trong bồn thì sao? Ai rửa ạ?''

Kuon lấy chiếc tạp dề vải từ móc treo trong bếp, ném nó nhẹ lên lưng ghế: ``Hai người cùng dọn dẹp bếp tối nay. Không phải để phân định ai thắng ai thua, không phải để phạt ai. Chỉ là vì bây giờ cả hai đều ở đây, cùng nhau làm thì nó nhanh thôi. Cứ đẩy qua đẩy lại thêm lại tốn thời gian chứ làm gì.''

Kaigou định nói gì đó, muốn phản đối hay đề xuất gì đó, nhưng Kuon giơ tay ngăn lại ngay. Anh nói: ``Không cần phải phản biện gì cả. Anh không phạt ai cả. Đây chỉ là việc nhà mà mọi người cùng làm thôi.''

Ryuuhou đứng dậy từ ghế, xắn nhẹ tay áo lên đến khuỷu tay, sẵn sàng làm việc. Em hỏi Kaigou: ``Vậy em rửa bát, chị lau khô chúng rồi xếp vào giá nhé? Được không?''

Kaigou nhìn em, mỉm cười: ``Được. Nhưng nhớ tráng kỹ bọt xà phòng ra khỏi đĩa đấy, đừng để còn sót lại.''

``Em biết rửa bát mà chị, đừng coi em như trẻ con,'' Ryuuhou cười phỏng, cuối cùng cũng hết đi sự căng thẳng từ nãy giờ.

``Vậy thì tốt, cứ thế mà làm.'' Kaigou cũng cười theo.

Kuon đặt chiếc tạp dề lên lưng ghế, không mặc vào, anh nói: ``Anh sẽ rửa chiếc nồi lớn ở trong bồn, cái nồi dùng để nấu lẩu hôm qua. Hai người lo phần bát đĩa là được.''

Sự căng thẳng cuối cùng cũng tan biến hoàn toàn, thay vào đó là không khí ấm áp của một gia đình.

\ct

Kaigou xắn cao cả hai tay áo lên khuỷu tay, bước đến bồn rửa chén chính, còn Kuon đặt chiếc nồi lớn đã dùng nấu lẩu hôm qua vào bồn phụ bên cạnh để rửa riêng. Ryuuhou xoay vòi nước, dùng đầu ngón tay thử nhiệt độ ấm vừa phải trước khi vặn nước to ra. Naomi nhanh nhẹn dời chồng khăn ăn đã xếp gọn sang góc bàn, nhường hết không gian giữa bếp cho mọi người làm việc.

Tôi lục trong tủ chén lấy miếng bọt biển mới còn nguyên, chưa dùng lần nào, rồi đưa cho chị: ``Onee-sama, miếng này mềm hơn cái cũ, lau bát không bị xước ạ.''

Chị nhận lấy miếng bọt biển, mỉm cười với tôi: ``Cảm ơn Suzu nhé.''

Ryuuhou cầm chai nước rửa chén từ trên kệ, chị Kaigou vừa vắt khăn lau bát vừa nhắc: ``Chỉ cần bóp ít thôi, nhiều bọt lại tốn nước tráng.''

Em ấn vòi chai, cố tình chỉ lấy một giọt, nhưng tay em trượt nhẹ, toàn bộ lượng xà phòng trong nhát bóp ấy tuôn ra nhiều hơn gấp mười lần. Khi nước đánh vào, bọt trắng phồng lên ngập cả miệng bồn. 

Kaigou ngẩng nhìn tay em, không giấu được sự bất lực: ``Đó không phải là một giọt đâu Ryuu.''

``Tay em trượt thôi mà!'' Ryuuhou cãi lý, vẫy vẫy đôi tay trắng trẻo của mình.

``Chai có vòi nhấn mà, ấn nhẹ thôi là ra được.''

``Em biết mà, hôm nay tay hơi trơn thôi.''

Từ bồn phụ, Kuon vừa cọ chiếc nồi vừa bình thản xen vào: ``Hai người có thể tranh luận xem cái `một giọt' của em to cỡ nào sau khi rửa xong tất cả bát đĩa cũng được. Bây giờ cứ làm việc trước đi.''

Kaigou cắn nhẹ môi, không đáp lại gì thêm, chỉ lắc đầu cười. Ryuuhou cũng cúi đầu bắt đầu cọ từng chiếc bát, tiếng nước chảy đều đều lấp đầy khoảng im lặng ngắn ngủi. Chiếc bát đầu tiên được tráng sạch vết bẩn, Kaigou lau khô bằng khăn bông rồi đặt gọn vào giá treo. Chiếc thứ hai, thứ ba cũng lần lượt được làm sạch, không ai nói thêm lời nào trong gần một phút.

Kson và Delutaya đã đặt túi bắp rang xuống ghế, mắt vẫn dán vào bếp. Naomi ngồi trên ghế đẩu nhỏ, chăm chú nhìn hai chị, tôi không thấy chút tức giận hay căng thẳng nào trên mặt họ nữa, chỉ còn sự tập trung nhẹ nhàng vào công việc đơn giản trước mắt.

\ct

Ryuuhou rửa bát nhanh hơn tốc độ chị Kaigou lau khô, chị phải cầm khăn lau thật kỹ, lật cả mặt sau đĩa để kiểm tra không còn vết dầu mỡ sót lại. Em vừa cọ xong chiếc đĩa tiếp theo, đưa sang cho chị nhưng chưa nhận ngay, chỉ chỉ vào mép đĩa: ``Còn sót bọt ở đây này.''

Ryuuhou ngước nhìn, quả nhiên có vệt bọt trắng nhỏ dính ở cạnh đĩa. Em liền xoay vòi nước lại, đưa đĩa dưới vòi tráng thêm lần nữa. ``Được rồi này.''

Lần này Kaigou mới nhận lấy, lau khô và đặt vào giá. Kaigou nói nhỏ: ``Cảm ơn em.''

Em hơi ngạc nhiên, ngẩng lên nhìn chị một cái rồi cũng mỉm cười: ``Không có gì.''

Đúng lúc đó, một chiếc thìa sắt bị tuột tay Ryuuhou, rơi xuống đáy bồn tạo tiếng leng keng rõ ràng. Kaigou cúi xuống nhặt nó lên, Ryuuhou vội tránh sang bên để nhường chỗ, hai người lại gần nhau đến mức vai gần như đụng vào nhau.

``Em đứng qua bên trái đi,'' Kaigou nói, vừa lùi lại một chút.

``Chị đang đứng bên trái rồi mà,'' Ryuuhou chỉ vị trí chị đang đứng.

``Vậy em sang bên phải đi.''

``Bên phải có giá bát rồi, đứng vào đó làm sao đặt bát?''

Anh Kuon ở bên cạnh ngẩng nhìn cảnh hai người đang xoay sở khó xử, chỉ nói một câu: ``Hai người đổi vị trí cho nhau là được mà.''

Nghe lời anh, hai người cùng bước sang phía trái một nhịp, rồi nhận ra vẫn chồng lấn, lại cùng bước về phía phải. Cảnh tượng vụng về đó khiến Naomi bật cười khúc khích đầu tiên. Kaigou và Ryuuhou cùng dừng lại, nhìn em nhỏ.

``Em cười gì thế?'' Ryuuhou hỏi.

Naomi đưa hai tay lên bắt chước động tác hai chị vừa bước qua bước lại, nói rõ ràng: ``Hai chị giống đang nhảy đôi ạ!''

Cả Delutaya phía xa cũng không nhịn được cười, Kson thì bật cười lớn, tiếng cười lan tỏa khắp phòng, không khí căng thẳng cuối cùng trong bếp tan biến hết, thay bằng sự ấm áp của một gia đình.

\ct

Khi chồng bát đĩa dần vơi đi nửa số lượng, trong bồn rửa chỉ còn lại vài chiếc đĩa và chiếc bát bị bỏ lại từ tối, tôi đứng ở góc bếp vừa xếp khăn ăn vừa quan sát hai chị. 

Lúc này chị Kaigou, vừa lau khô xong chiếc đĩa gốm sứ màu trắng, bỗng dừng tay và quay sang hỏi em gái mình: ``Hôm thứ Ba em mang chồng khăn tắm lên máy sấy tầng sáu à?''

Ryuuhou, lúc này đang cọ bọt xà phòng lên chiếc bát cuối cùng, gật đầu nhanh chóng, nước còn nhỏ giọt từ cổ tay em xuống sàn nhà chống trượt: ``Dạ. Máy giặt vừa chạy xong, khăn ướt vẫn nằm nguyên trong giỏ nhựa chờ mang lên. Em thấy hôm đó chị bận đi làm về muộn, nên tự ý mang lên luôn.''

``Chị\dot\ cũng không hề hay biết việc đó đâu.'' Kaigou chần chừ một chút, đặt chiếc đĩa vừa lau khô vào giá treo, tiếng đĩa va nhẹ vào nhau tạo ra âm thanh trong trẻo.

``Em quên nhắn vào nhóm thông báo,'' Ryuuhou thừa nhận, vội vàng tráng chiếc bát dưới vòi nước lạnh để gột sạch hết bọt xà phòng.

Kaigou xoay người, nhìn em gái với ánh mắt dịu đi hẳn so với lúc đầu buổi tối: ``Chị còn thấy chồng khăn được gấp gọn trên giá phòng tắm, còn tưởng là Suzu làm chứ.''

Ryuuhou khẽ nhún vai, nhưng góc môi em hơi nâng lên một chút, như thể nhẹ nhõm vì cuối cùng cũng có người biết em đã làm việc đó.

Em tráng xong chiếc bát cuối cùng, đặt nó vào chồng bát sạch chờ lau khô, rồi gật đầu: ``Lần sau em sẽ nhắn vào nhóm ngay khi làm thêm việc gì ngoài lịch để mọi người cùng biết ạ.''

Kaigou cười nhẹ, lau khô chiếc bát vừa được Ryuuhou rửa sạch: ``Còn chị thì hứa sẽ không viết dấu sao đỏ bên cạnh tên em trước khi hỏi em rõ lý do nữa.''

``Thế thì chị có thể viết dấu sao xanh thay thế được không?'' Ryuuhou ngay lập tức gợi ý, còn cố tình nháy mắt với chị gái làm Kaigou phải lắc đầu cười. 

``Chị không có ý định đổi mực đâu.''

``Vậy thì dấu sao vàng? Dấu nào cũng được, miễn là không phải màu đỏ xui xẻo đó!'' cô em gái tiếp tục đùa, lấy khăn giấy lau tay khô.

``Ryuu. Chị không đùa đâu đấy nhé,'' chị Kaigou gọi tên em, giọng giả vờ nghiêm nghị nhưng không che giấu được nụ cười.

``Em chỉ gợi ý thôi mà!''

Ngay lúc đó, chiếc đồng hồ thông minh trên bàn ăn bỗng kêu bíp một tiếng, Game lại tiếp tục công việc của mình: ``Tôi có thể thêm cột `đã xác nhận' vào lịch phân công việc nhà chung. Mọi thay đổi đều cần có xác nhận của cả hai bên trước khi cập nhật.''

``Thêm đi,'' chị lập tức đáp, không chần chừ gì nữa.

Ryuuhou cũng gật đầu lia lịa, thêm vào yêu cầu của mình: ``Và thêm cột `đổi ca đã được cả hai đồng ý' nữa nhé. Để tránh tình trạng người nhận không biết mình phải làm việc đó.''

``Tôi sẽ cập nhật mẫu lịch ngay lập tức, yêu cầu của hai người đã được ghi nhận.'' Game trả lời đều đều như mọi khi, màn hình đồng hồ nhấp nháy xanh lá báo hiệu đang xử lý.

Kuon lau khô chiếc nồi lớn, đặt nó lên giá treo trên cao khỏi tầm tay các em nhỏ, rồi quay sang nhìn hai chị em, giọng anh mang chút trêu chọc: ``Không cần thêm cột về cảm xúc hay tâm trạng hiện tại vào lịch chứ? Để ai cũng biết hôm nay ai đang vui hay buồn mà tránh?''

Game trả lời ngay lập tức, không có chút do dự: ``Tôi không có dữ liệu đáng tin cậy để theo dõi và cập nhật cột đó. Các cảm xúc thay đổi quá nhanh để tôi có thể ghi nhận chính xác.''

Ryuuhou khẽ cười thành tiếng, tiến lại gần bàn ăn để nhìn màn hình đồng hồ đang cập nhật: ``Vậy thì để bọn tôi tự nói với nhau về cảm xúc của mình thôi. Công nghệ đâu thể thay thế được việc hai người ngồi lại chia sẻ.''

\ct

Bồn rửa cuối cùng cũng trống trơn, không còn một chiếc bát đĩa nào sót lại. Toàn bộ bọt xà phòng đã được xả sạch, vòi nước được đóng kín. Mặt bàn ăn vốn có vài vệt nước cũng được lau khô hoàn toàn, sáng bóng dưới ánh đèn vàng ấm của phòng khách. Chiếc khăn lau bông đã được vắt kiệt nước, treo lại đúng móc sắt sau cửa bếp như ban đầu.

Kaigou kiểm tra lại vòi nước chính một lần nữa, đảm bảo không có hiện tượng rò rỉ nhỏ nào. Ryuuhou cất chai nước rửa chén trở lại kệ trên cao, xếp nó thẳng hàng với các chai tẩy rửa khác. Kuon rửa sạch tay dưới vòi nước phụ, lau khô bằng khăn giấy rồi tắt đèn bếp phụ, để lại chỉ có ánh sáng từ đèn bếp chính chiếu rọi toàn bộ căn phòng.

Trên bàn phòng khách, Game đã hoàn thành việc cập nhật và hiển thị mẫu lịch phân công việc nhà mới lên màn hình lớn, rõ ràng hơn bao giờ hết. Các cột được sắp xếp tuần tự:

\begin{itemize}
    \item Tên công việc cần hoàn thành;
    \item Người được phân công nhận việc;
    \item Thời hạn hoàn thành;
    \item Trạng thái đã được xác nhận hay chưa;
    \item Và một ô ghi chú nhỏ ở cuối cùng, để điền thông tin khi có nhu cầu đổi ca giữa các thành viên.
\end{itemize}

Kaigou tiến lại gần bàn, đọc qua từng dòng trên màn hình, ngón tay chỉ nhẹ vào màn hình để kiểm tra các mục: ``Được rồi. Như thế này rõ hơn mẫu cũ rất nhiều. Không còn lý do để có hiểu lầm nữa.''

Ryuuhou cũng đứng cạnh chị, nhìn vào dòng tên của mình được liệt kê trong lịch ngày mai: ``Tôi nghĩ từ giờ sẽ không phải dùng đến bảng công trạng của mình nữa. Mọi thứ đã được ghi nhận rõ ràng trên lịch chung rồi.''

Chị khép cuốn sổ màu xanh lam cũ mòn của mình lại, đặt nó nhẹ nhàng lên góc bàn, không còn để nó trải rộng như một bản truy tố nữa: ``Chị vẫn sẽ giữ cuốn sổ này.''

Ryuuhou quay sang nhìn chị, mày hơi nhíu lại trong thắc mắc: ``Để ghi những việc chưa được phân công vào đó à? Những việc phát sinh ngoài lịch?''

``Để nhớ những gì chị đã học được từ hôm nay.'' Kaigou nói chậm rãi, nhìn vào bìa cuốn sổ có vài vết ố cũ. ``Nhớ rằng không phải mọi thứ đều có thể ghi vào sổ sách hay bảng biểu. Đôi khi cần phải ngồi lại nói chuyện mới hiểu được người kia đang nghĩ gì.''

Ryuuhou hơi khựng lại một chút, như thể không ngờ chị sẽ nói vậy. Rồi em cũng gật đầu đồng tình, tay chạm nhẹ vào chiếc bảng kẹp giấy công trạng của mình đang nằm trên ghế: ``Vậy thì chắc em cũng giữ bảng của mình.''

Chị Kaigou nhướng mày, trêu chọc em gái: ``Em vừa nói là không cần dùng nữa mà? Sao bây giờ lại giữ?''

``Em sẽ dùng nó để ghi những ý tưởng nấu ăn mới tìm thấy trên mạng.'' Ryuuhou giải thích, cầm chiếc bảng kẹp lên, vuốt nhẹ dòng chữ ``bảng công trạng'' viết tay ở tiêu đề. ``Ghi công thức nấu ăn, món mới muốn thử làm cho cả nhà.''

``Bảng công trạng biến thành sổ công thức nấu ăn à? Thế có đáng để chuyển đổi không?'' Kaigou cười.

``Em chưa nghĩ ra tên mới cho nó đâu. Để từ từ nghĩ cũng được.'' Ryuuhou nhún vai, đặt chiếc bảng vào túi sách của mình.

Kuon đứng lên từ ghế, cầm túi công việc của mình chuẩn bị lên phòng, rồi quay lại nhắc nhở hai chị em lần cuối: ``Miễn là lần sau không dùng cái bảng đó hay cuốn sổ đó để tranh luận ai đã lau bàn nhiều hơn, ai đã rửa bát nhiều hơn là được. Đừng để chuyện nhỏ nhặt như vậy lại gây căng thẳng nữa.''

Ryuuhou nhìn chị Kaigou, rồi quay sang đáp với anh Kuon, giọng đầy tinh nghịch: ``Em không hứa đâu. Nếu có tranh cãi thì em vẫn sẽ dùng bảng của mình làm bằng chứng đấy nha\~{}''

Kaigou thở dài, nhưng lần này chị hoàn toàn đang cười, không có chút bực bội nào trong giọng nói. Cả gian bếp lại tràn ngập tiếng cười nhẹ nhàng, ấm áp của một gia đình, cuối cùng mọi căng thẳng từ đầu buổi tối đã tan biến hoàn toàn.

\ct

Kson vươn tay nhấc túi bắp rang đã gần cạn lên, miệng vẫn còn mấp máy vài hạt bơ còn sót lại. Chị ngước mắt nhìn quanh căn bếp vừa mới sạch bóng, rồi hỏi với giọng đầy tiếc nuối: ``Vậy là chương trình tối nay chính thức kết thúc rồi à? Em còn chuẩn bị sẵn bài bình luận cuối cùng cho phần trao giải nữa cơ.''

Kaigou vừa gấp chiếc khăn lau bát lại, ngẩng lên nhìn Kson với ánh mắt khó hiểu: ``Em nói cái gì thế? Có chương trình gì để kết thúc đâu.''

``Tiếc thật đấy,'' Kson thở dài giả tạo, vỗ nhẹ lên chiếc sofa nỉ. ``Em đã dành cả buổi để chuẩn bị phần nhận xét công bằng cho cả hai bên mà.''

Delutaya kéo nhẹ góc tay áo của Kson, giọng nhỏ nhắn xen vào:
``Chị cứ giữ lại bài nhận xét đó đi. Còn nhiều buổi `biểu diễn' khác sau này mà.''

Vừa lúc đó, Naomi chạy nhảy đến bên Kaigou, đôi bàn tay bé xíu giơ cao một tờ giấy vẽ đang cẩn thận giữ ở hai góc. Cô bé nói với giọng hào hứng, má còn hơi ửng hồng vì vừa chạy qua chạy lại:
``Kaigou-oneechan, Ryuuhou-oneechan, nhìn này! Naomi đã vẽ hai chị rồi này!''

Tôi đứng nép ở góc bếp, nhìn Naomi dang rộng tờ giấy ra. Trên đó là hai hình người chibi tròn trịa đứng đối mặt nhau ở hai bên một chiếc bồn rửa chén khổng lồ, cao gần bằng cả hai người đứng cạnh. Một người cầm chặt cuốn sổ màu xanh lam trên tay, lông mày hơi cau lại trông rất nghiêm nghị. Người còn lại thì gồng mình giữ chiếc bảng kẹp giấy to hơn cả người mình, miệng cười rộng hết cỡ. Chính giữa chiếc bồn rửa, Game được vẽ thành một chiếc đồng hồ tròn nhỏ có đôi mắt to tròn long lanh, và Naomi còn tô một vệt nước xanh lè chạy dài khắp mặt bồn, như thể có cả đống bọt xà phòng vừa tràn ra.

Chị Kaigou chăm chú nhìn bức tranh một lúc, rồi chỉ vào chiếc bồn rửa khổng lồ hỏi: ``Sao em vẽ cái bồn rửa to thế kia? Nó còn to hơn cả hai chị cộng lại nữa.''

Naomi đáp ngay, không chần chừ gì: ``Vì lúc đó cái bồn đó đúng là một chiến trường mà! Hai chị đứng đối đầu nhau ở hai bên, nó giống hệt trong phim hoạt hình Naomi xem ạ!''

Ryuuhou đứng bên cạnh bật cười khúc khích, tay che miệng để khỏi cười quá to. Kaigou ban đầu còn cố nén cười, giữ cho khuôn mặt trông nghiêm nghị một chút, nhưng cuối cùng cũng không chịu được mà bật cười lớn, đôi vai run lên nhè nhẹ.

``Em vẽ chị trông dữ dằn quá đấy,'' chị chỉ vào hình người cầm sổ có cặp lông mày cau lại. ``Chị có đáng sợ thế này đâu.''

``Naomi vẽ đúng lúc chị đang tức giận mà!'' cô bé nhún nhảy trên chân. ``Lúc đó chị trông đúng là vậy ạ!''

Ryuuhou cũng nghiêng đầu nhìn bức tranh, rồi chỉ vào hình người còn lại hỏi với giọng tò mò: ``Còn chị kia thì trông thế nào ạ? Có dữ không?''

Naomi liền chỉ ngón tay bé xíu vào hình cầm bảng kẹp: ``Ryuuhou-oneechan đang cười mà! Nhìn miệng cười rộng thế này.''

``Đúng rồi! Chị lúc đó cũng đang cười thật mà!'' Ryuuhou vỗ tay vui vẻ.

Naomi cẩn thận đưa bức tranh cho hai chị. Kaigou đỡ lấy góc trái của tờ giấy, Ryuuhou thì giữ góc phải. Không ai giành giật hay tranh cãi gì nữa, cả hai cùng nâng bức tranh lên, giữ nó thẳng tắp ở giữa để ngắm nghía tác phẩm của em nhỏ. Khoảnh khắc đó thật ấm áp, tất cả sự căng thẳng từ mấy tiếng trước đã tan biến hoàn toàn, thay vào đó là niềm vui đơn giản của một gia đình.

\ct

Khi Kson và Delutaya mang theo túi bắp rang còn sót vài hạt xuống nhà, tầng sáu dần yên tĩnh trở lại, chỉ còn tiếng đồng hồ tích tắc đều đều trong góc phòng. 

Tôi là người cuối cùng còn ở lại bàn ăn, thu dọn nốt chồng khăn ăn mà Naomi đã xếp từ nãy giờ, xếp chúng gọn gàng vào ngăn kéo tủ bếp. Naomi thì cất hộp bút màu của mình vào túi xách nhỏ, đậy nắp hộp thật cẩn thận trước khi chạy lên cầu thang.

Chị Kaigou mở ngăn kéo dưới bàn ăn, đặt cuốn sổ màu xanh lam cũ mòn của mình vào trong, đóng ngăn kéo lại thật nhẹ nhàng như thể đang cất giữ một kỷ vật quan trọng. Ryuuhou thì lấy chiếc bảng kẹp công trạng của mình, kẹp nó chặt vào sau cuốn lịch học dày của mình, để nó không bị nhàu hay gãy góc. Không ai xé đi hay vứt bỏ những thứ đã từng là ``bằng chứng'' trong cuộc tranh cãi vừa rồi.

Cả hai người dường như đều hiểu rằng những tờ giấy đó, những danh sách và ghi chép đó không hề sai. Chúng chỉ đơn thuần là những con số, những ghi chú vô tri. Chúng không thể thay thế được một cuộc trò chuyện chân thành, những lời nói thật lòng mà hai chị em vừa trải qua.

Trước khi bước lên cầu thang lên phòng, Ryuuhou đột nhiên dừng lại ở cửa bếp, quay lại gọi chị gái mình: ``Chị Kaigou này.''

Kaigou đang lau lại mặt bàn một lần nữa, nghe tiếng gọi thì quay người lại: ``Cái gì thế em?''

``Ngày mai em sẽ rửa bát sau bữa tối,'' em nói rõ ràng, gật đầu như thể đã sắp xếp mọi thứ trong đầu. ``Em sẽ đợi mọi người ăn xong rồi rửa hết.''

``Chị sẽ đánh dấu vào lịch chung ngay bây giờ,'' Kaigou mỉm cười, gật đầu đồng ý.

``Nhưng nếu hôm đó em có làm thêm món gì khác, hay có việc bận đột xuất, em sẽ nhắn tin vào nhóm ngay.''

``Và chị sẽ trả lời xác nhận ngay lập tức. Không bao giờ để em phải chờ đợi nữa,'' Kaigou nói chắc chắn.

Ryuuhou cười lớn, gật đầu lia lịa: ``Được rồi ạ!''

Kaigou cũng gật đầu, giọng nhẹ nhàng: ``Ừ, được rồi.''

Ryuuhou quay người bước xuống cầu thang, tiếng bước chân nhỏ dần dần lên tầng năm. Chị Kaigou đứng yên ở cửa bếp, nhìn theo bóng em gái mình biến mất ở cầu thang, rồi lại cầm chiếc khăn lau bông lên. Tôi biết mặt bàn đã sạch bóng từ lâu, không cần phải lau lại lần nào nữa. Nhưng chị vẫn muốn kiểm tra một lần cuối, như thể đang tự nhủ rằng mọi thứ đã thực sự ổn thỏa, không còn gì phải lo lắng nữa.

%==========================================TRUYỆN PHỤ=========================================
\omk

Sáng hôm sau, khi mọi người vừa thức dậy, tin nhắn đầu tiên xuất hiện trong nhóm chat gia đình là từ Game. Thông báo được gửi đi cùng một tệp đính kèm lịch phân công việc nhà mới, nội dung rõ ràng: ``Mẫu phân công việc nhà đã được cập nhật lên hệ thống. Mọi thay đổi, yêu cầu đổi ca hay thêm việc đều cần có xác nhận bằng tin nhắn từ cả hai bên trước khi được cập nhật vào lịch chung.''

Chưa đầy một phút sau, Ryuuhou gửi biểu tượng ngón tay cái lớn, kèm theo một sticker cười vui vẻ. Kaigou gửi một dấu tick xanh xác nhận đã đọc và đồng ý. Naomi thì gửi một hình trái tim đỏ rực, đáng yêu hết mực. 

Kson gửi bất ngờ một bức ảnh chụp chồng bát sạch bóng trong bếp, kèm theo một dấu chấm hỏi. Delutaya liền nhắn lại hỏi ai đã rửa chúng sạch thế. Chị ấy trả lời ngay: ``Không phải chị.''

Game lập tức nhắn lại: ``Bức ảnh được cung cấp không đủ dữ liệu để xác định người đã rửa chén. Không thể cập nhật vào lịch phân công.''

Tôi mở điện thoại ra đọc tin nhắn, thì thấy chị Kaigou nhắn vào nhóm, dòng chữ đơn giản nhưng ấm áp: ``Không cần xác định đâu. Ai rửa thì người đó đã làm. Không cần phải ghi vào lịch làm gì.''

Ryuuhou, lúc đó đang trên đường đến trường, đọc được tin nhắn đó liền gửi lại biểu tượng ngón tay cái một lần nữa. Tôi có thể thấya em đang mỉm cười trên đường đi, gió thổi bay mái tóc của mình.

\ct

Đến buổi tối hôm đó, sau khi bữa ăn đã kết thúc và mọi người đang ngồi thư giãn trong phòng khách, Naomi lúi húi mang bức tranh chibi mà em đã vẽ ra, hai tay cầm chắc hai góc tờ giấy để không bị nhàu nát. Em bước nhẹ ra giữa phòng, đặt bức tranh cẩn thận lên chiếc bàn ăn lớn, ngay cạnh chiếc đồng hồ thông minh của Game vẫn đang nhấp nháy nhẹ nhàng. 

``Naomi muốn tất cả mọi người cùng ký tên vào bức tranh này ạ!'' cô bé reo lên, đôi mắt tròn long lanh nhìn quanh từng người trong phòng.

Nghe lời em nhỏ, chị Kaigou là người đầu tiên đứng dậy, cầm cây bút bi đen từ trên bàn rồi viết chữ ký của mình ngay dưới hình người cầm cuốn sổ xanh lam mà Naomi đã vẽ. Tiếp theo, Ryuuhou tranh thủ cầm lấy cây bút trước khi người khác giành mất, ký tên mình rực rỡ bên cạnh hình người đang ôm chiếc bảng kẹp công trạng to lớn.

Tôi cũng bước đến, lấy một cây bút màu hồng từ hộp bút của Naomi, vẽ thêm một trái tim nhỏ xinh ở góc dưới cùng của tờ giấy, nơi còn trống trải. Anh Kuon thì cười lớn, lấy bút màu lam vẽ một chiếc bát sạch bóng bên cạnh cái bồn rửa khổng lồ trong tranh, như muốn thêm vào ``thành quả'' của buổi tối hôm trước.

Kson không chịu kém cạnh, viết dòng chữ to đùng ngay bên cạnh bức vẽ của mình: ``Khán giả khách mời đặc biệt.'' Còn Delutaya thì lấy bút màu vàng, vẽ thêm một túi bắp rang nhỏ trong tay nhân vật Kson trong tranh, đúng như những gì đã xảy ra tối qua.

Bỗng nhiên, chiếc đồng hồ của Game phát ra tiếng bíp nhỏ, giọng điện tử đều đều vang lên: ``Tôi có nên ký tên vào đây không? Tôi không có tay để cầm bút đâu.'' 

Nghe vậy, Naomi lập tức nhặt cây bút màu xanh lá cây đưa về phía chiếc đồng hồ, như thể Game có thể nhận được, rồi nói to: ``Naomi sẽ cầm bút giúp Game-san ạ!'' Em đặt đầu bút lên góc trống còn lại của tờ giấy, rồi viết dòng chữ ``Game-san'' thật ngay ngắn, nét chữ dù chưa thành thạo nhưng rất rõ ràng. 

Sau khi tất cả mọi người đã ``ký tên'', cô bé lấy cuộn băng dính trong suốt từ ngăn kéo tủ, dán bức tranh lên cánh cửa tủ lạnh lớn ở giữa bếp, nơi mọi người có thể nhìn thấy nó mỗi khi đi qua.

Chị Kaigou đứng nhìn bức tranh trên tủ lạnh một lúc, rồi giả vờ thở dài nói với em gái mình: ``Lần sau đừng vẽ cái bồn rửa to như cái hồ thế này nữa nhé, trông nó còn to hơn cả hai chị em cộng lại.'' 

Ryuuhou lập tức đáp lại, tay chống hông trông rất thuyết phục: ``Nhưng nó là nhân vật chính của bức tranh mà! Không có bồn rửa thì câu chuyện không có gì để nói nữa!''

Nghe vậy, Kaigou cầm chiếc khăn lau bàn vừa xếp gọn lên, giả vờ nổi giận đuổi theo Ryuuhou quanh phòng. Cô em gái cười lớn, chạy vòng quanh chiếc bàn ăn để trốn tránh chị, trong khi Naomi đứng cạnh tủ lạnh cười vang lên, tiếng cười của em nhỏ lan tỏa khắp căn phòng.  Ngay cả Game cũng nhắc nhở bằng giọng đều đều: ``Đề nghị hai người không chạy gần khu vực bồn rửa, có thể gặp nguy hiểm nếu trượt ngã.'' 

Anh Kuon thì ngồi trên ghế sofa, thở dài như thể đang lo lắng cho hai người chị em không bao giờ lớn được, nhưng nụ cười không bao giờ tắt trên môi anh, làm cho không khí trong căn nhà càng thêm ấm áp.

Tối qua, không ai thắng được trận ``đấu khẩu'' về việc ai làm nhiều ai làm ít. Không ai có thể khoe mình là người đúng hoàn toàn. Nhưng sáng nay, bồn rửa bát đã trống trơn, sạch bóng. Và lần sau, nếu có bất kỳ tranh cãi nào nảy ra, có lẽ sẽ không cần ai phải lập hồ sơ, phải viết danh sách hay chuẩn bị bằng chứng gì nữa để được lắng nghe. Chỉ cần hai người ngồi lại, nói chuyện chân thành với nhau, là mọi thứ sẽ được giải quyết.

\end{document}